\documentclass[10pt]{extarticle}
\usepackage[a4paper,margin=1.5cm]{geometry}
\usepackage[T1]{fontenc}
\usepackage{amsmath,amssymb}
\usepackage{graphicx}
\usepackage[export]{adjustbox}
\usepackage{multicol}
\usepackage{enumitem}
\usepackage{booktabs,array,tabularx}
\setlength{\columnsep}{0.7cm}
\setlength{\parindent}{0pt}
\setlength{\parskip}{3pt}
\pagestyle{empty}
\makeatletter
% figure registry (figdefs.tex is generated by the extract step)
\newcommand{\qfigdef}[4]{\@namedef{qf@file@#1}{#2}\@namedef{qf@w@#1}{#3}\@namedef{qf@wide@#1}{#4}}
\newcommand{\qf@check}[1]{\@ifundefined{qf@file@#1}{\PackageError{score2}{Unknown figure id #1}{}}{}}
\newcommand{\qf@img}[4][]{\includegraphics[width=#2\dimexpr\@nameuse{qf@w@#4}pt\relax,max width=#3,max height=0.45\textheight,#1]{\@nameuse{qf@file@#4}}}
% Full-width questions: finalize wraps every question that contains a wide figure in \qfullbegin...\qfullend,
% so the WHOLE question (stem, figures, options) is set across the page in reading order, then the two columns
% resume. The question's head - its number, stem and first figure (\qfullsplit marks the end of the head) -
% is kept on one page, so a number is never left at a page foot with its figure overleaf; the rest may break
% across pages normally, so no big gaps are left. Without a \qfullsplit, the whole question is kept together
% if it is under 45% of a page. Never write these yourself.
\newif\ifqf@full
\newif\ifqf@boxed
\newsavebox{\qf@box}
\newcommand{\qf@open}{\global\qf@fulltrue\par\global\qf@boxedtrue
  \setbox\qf@box\vbox\bgroup\hsize\textwidth\linewidth\textwidth}
\newcommand{\qf@emit}[1]{\par\egroup\global\qf@boxedfalse
  \ifdim\dimexpr\ht\qf@box+\dp\qf@box\relax<#1\textheight\noindent\usebox{\qf@box}\par\else\unvbox\qf@box\fi}
\newcommand{\qfullsplit}{\ifqf@boxed\qf@emit{0.85}\fi}
\newcommand{\qf@close}{\ifqf@boxed\qf@emit{0.45}\fi\global\qf@fullfalse}
\newcommand{\qfullbegin}{\end{multicols}\qf@open}
\newcommand{\qfullend}{\qf@close\begin{multicols}{2}}
\newcommand{\qfullstart}{\qf@open}
\newcommand{\qfulldone}{\qf@close}
% \qfig[scale]{id}: block figure, centred, directly where it is written (size set by the tool).
% Fallback if a wide figure is not inside a full-width question: break out of the columns in place - no floats.
\newcommand{\qfig}[2][1]{\qf@check{#2}%
  \ifqf@full\qf@blk{#1}{#2}\else
  \ifnum\@nameuse{qf@wide@#2}=1 \expandafter\@firstoftwo\else\expandafter\@secondoftwo\fi
  {\end{multicols}\begin{center}\qf@img{#1}{\textwidth}{#2}\end{center}\begin{multicols}{2}}%
  {\qf@blk{#1}{#2}}\fi}
\newcommand{\qf@blk}[2]{\par\smallskip{\centering\qf@img{#1}{\linewidth}{#2}\par}\smallskip}
% \qfigi[scale]{id}: inline figure (inside option lists / tables / minipages); never breaks columns.
\newcommand{\qfigi}[2][1]{\qf@check{#2}\qf@img[valign=c]{#1}{\linewidth}{#2}}
\makeatother
% structure macros - lint reads these to check completeness and figure placement
\newcommand{\papersec}[1]{\par\bigskip{\centering\Large\bfseries PAPER-#1\par}\medskip}
\newcommand{\partsec}[2]{\par\medskip\noindent{\large\bfseries PART-#1: #2}\par\smallskip}
\newcommand{\sect}[1]{\par\smallskip\noindent\textbf{#1}\par}
\newcommand{\qn}[1]{\par\smallskip\noindent\textbf{#1.}~}
\newcommand{\paraq}[2]{\par\smallskip\noindent\textbf{Paragraph for Questions #1 and #2}\par}
\newenvironment{qoptions}{\begin{itemize}[leftmargin=1.8em,itemsep=1pt,topsep=2pt,label={}]}{\end{itemize}}
\newcommand{\opt}[2]{\item[(#1)] #2}
% figure registry (generated by score2_tools.py finalize)
\qfigdef{p1_phy_s1_q01_1}{TAAS - Test 02 (30 Jul 2023)-img/p1_phy_s1_q01_1.png}{126.9}{0}
\qfigdef{p1_phy_s1_q02_1}{TAAS - Test 02 (30 Jul 2023)-img/p1_phy_s1_q02_1.png}{108.3}{0}
\qfigdef{p1_phy_s1_q03_1}{TAAS - Test 02 (30 Jul 2023)-img/p1_phy_s1_q03_1.png}{107.5}{0}
\qfigdef{p1_phy_s2_q06_1}{TAAS - Test 02 (30 Jul 2023)-img/p1_phy_s2_q06_1.png}{71.1}{0}
\qfigdef{p1_phy_s2_q06_A}{TAAS - Test 02 (30 Jul 2023)-img/p1_phy_s2_q06_A.png}{68.5}{0}
\qfigdef{p1_phy_s2_q06_B}{TAAS - Test 02 (30 Jul 2023)-img/p1_phy_s2_q06_B.png}{68.9}{0}
\qfigdef{p1_phy_s2_q06_C}{TAAS - Test 02 (30 Jul 2023)-img/p1_phy_s2_q06_C.png}{68.5}{0}
\qfigdef{p1_phy_s2_q06_D}{TAAS - Test 02 (30 Jul 2023)-img/p1_phy_s2_q06_D.png}{68.9}{0}
\qfigdef{p1_phy_s2_q07_1}{TAAS - Test 02 (30 Jul 2023)-img/p1_phy_s2_q07_1.png}{84.6}{0}
\qfigdef{p1_phy_s3_q09_1}{TAAS - Test 02 (30 Jul 2023)-img/p1_phy_s3_q09_1.png}{245.2}{0}
\qfigdef{p1_phy_s3_q10_1}{TAAS - Test 02 (30 Jul 2023)-img/p1_phy_s3_q10_1.png}{137.0}{0}
\qfigdef{p1_phy_s3_q11_1}{TAAS - Test 02 (30 Jul 2023)-img/p1_phy_s3_q11_1.png}{46.3}{0}
\qfigdef{p1_phy_s3_q11_2}{TAAS - Test 02 (30 Jul 2023)-img/p1_phy_s3_q11_2.png}{35.0}{0}
\qfigdef{p1_phy_s3_q11_3}{TAAS - Test 02 (30 Jul 2023)-img/p1_phy_s3_q11_3.png}{38.7}{0}
\qfigdef{p1_phy_s3_q11_4}{TAAS - Test 02 (30 Jul 2023)-img/p1_phy_s3_q11_4.png}{33.0}{0}
\qfigdef{p1_phy_s3_q11_5}{TAAS - Test 02 (30 Jul 2023)-img/p1_phy_s3_q11_5.png}{41.3}{0}
\qfigdef{p1_phy_s4_q03_1}{TAAS - Test 02 (30 Jul 2023)-img/p1_phy_s4_q03_1.png}{76.4}{0}
\qfigdef{p1_phy_s4_q04_1}{TAAS - Test 02 (30 Jul 2023)-img/p1_phy_s4_q04_1.png}{97.9}{0}
\qfigdef{p1_phy_s4_q05_1}{TAAS - Test 02 (30 Jul 2023)-img/p1_phy_s4_q05_1.png}{92.5}{0}
\qfigdef{p1_phy_s4_q06_1}{TAAS - Test 02 (30 Jul 2023)-img/p1_phy_s4_q06_1.png}{79.4}{0}
\qfigdef{p1_chem_s1_q02_1}{TAAS - Test 02 (30 Jul 2023)-img/p1_chem_s1_q02_1.png}{32.0}{0}
\qfigdef{p1_chem_s1_q03_1}{TAAS - Test 02 (30 Jul 2023)-img/p1_chem_s1_q03_1.png}{203.5}{0}
\qfigdef{p1_chem_s2_q05_1}{TAAS - Test 02 (30 Jul 2023)-img/p1_chem_s2_q05_1.png}{41.8}{0}
\qfigdef{p1_chem_s3_q09_1}{TAAS - Test 02 (30 Jul 2023)-img/p1_chem_s3_q09_1.png}{193.4}{0}
\qfigdef{p1_chem_s3_q10_1}{TAAS - Test 02 (30 Jul 2023)-img/p1_chem_s3_q10_1.png}{84.0}{0}
\qfigdef{p1_chem_s3_q10_2}{TAAS - Test 02 (30 Jul 2023)-img/p1_chem_s3_q10_2.png}{68.1}{0}
\qfigdef{p1_chem_s3_q10_3}{TAAS - Test 02 (30 Jul 2023)-img/p1_chem_s3_q10_3.png}{72.2}{0}
\qfigdef{p1_chem_s3_q11_1}{TAAS - Test 02 (30 Jul 2023)-img/p1_chem_s3_q11_1.png}{107.8}{0}
\qfigdef{p1_chem_s4_q01_3}{TAAS - Test 02 (30 Jul 2023)-img/p1_chem_s4_q01_3.png}{28.4}{0}
\qfigdef{p1_chem_s4_q01_4}{TAAS - Test 02 (30 Jul 2023)-img/p1_chem_s4_q01_4.png}{41.1}{0}
\qfigdef{p1_chem_s4_q01_1}{TAAS - Test 02 (30 Jul 2023)-img/p1_chem_s4_q01_1.png}{32.0}{0}
\qfigdef{p1_chem_s4_q01_2}{TAAS - Test 02 (30 Jul 2023)-img/p1_chem_s4_q01_2.png}{34.6}{0}
\qfigdef{p1_chem_s4_q03_1}{TAAS - Test 02 (30 Jul 2023)-img/p1_chem_s4_q03_1.png}{143.5}{0}
\qfigdef{p1_chem_s4_q05_1}{TAAS - Test 02 (30 Jul 2023)-img/p1_chem_s4_q05_1.png}{208.2}{0}
\qfigdef{p2_phy_s1_q01_1}{TAAS - Test 02 (30 Jul 2023)-img/p2_phy_s1_q01_1.png}{145.2}{0}
\qfigdef{p2_phy_s1_q04_1}{TAAS - Test 02 (30 Jul 2023)-img/p2_phy_s1_q04_1.png}{148.0}{0}
\qfigdef{p2_phy_s2_q05_1}{TAAS - Test 02 (30 Jul 2023)-img/p2_phy_s2_q05_1.png}{108.3}{0}
\qfigdef{p2_phy_s2_q06_1}{TAAS - Test 02 (30 Jul 2023)-img/p2_phy_s2_q06_1.png}{67.1}{0}
\qfigdef{p2_phy_s2_q07_1}{TAAS - Test 02 (30 Jul 2023)-img/p2_phy_s2_q07_1.png}{58.5}{0}
\qfigdef{p2_phy_s3_para1_2_1}{TAAS - Test 02 (30 Jul 2023)-img/p2_phy_s3_para1_2_1.png}{129.3}{0}
\qfigdef{p2_phy_s3_para3_4_1}{TAAS - Test 02 (30 Jul 2023)-img/p2_phy_s3_para3_4_1.png}{105.3}{0}
\qfigdef{p2_phy_s4_q05_1}{TAAS - Test 02 (30 Jul 2023)-img/p2_phy_s4_q05_1.png}{108.3}{0}
\qfigdef{p2_phy_s4_q06_1}{TAAS - Test 02 (30 Jul 2023)-img/p2_phy_s4_q06_1.png}{86.2}{0}
\qfigdef{p2_phy_s4_q08_1}{TAAS - Test 02 (30 Jul 2023)-img/p2_phy_s4_q08_1.png}{161.1}{0}
\qfigdef{p2_phy_s4_q09_1}{TAAS - Test 02 (30 Jul 2023)-img/p2_phy_s4_q09_1.png}{211.3}{0}
\qfigdef{p2_chem_s1_q03_1}{TAAS - Test 02 (30 Jul 2023)-img/p2_chem_s1_q03_1.png}{61.0}{0}
\qfigdef{p2_chem_s3_para1_2_1}{TAAS - Test 02 (30 Jul 2023)-img/p2_chem_s3_para1_2_1.png}{282.6}{1}
\qfigdef{p2_chem_s4_q07_1}{TAAS - Test 02 (30 Jul 2023)-img/p2_chem_s4_q07_1.png}{136.4}{0}
\qfigdef{p2_chem_s4_q08_1}{TAAS - Test 02 (30 Jul 2023)-img/p2_chem_s4_q08_1.png}{129.4}{0}
\qfigdef{p2_chem_s4_q08_2}{TAAS - Test 02 (30 Jul 2023)-img/p2_chem_s4_q08_2.png}{129.4}{0}
\qfigdef{p2_chem_s4_q08_3}{TAAS - Test 02 (30 Jul 2023)-img/p2_chem_s4_q08_3.png}{128.1}{0}
\qfigdef{p2_chem_s4_q08_4}{TAAS - Test 02 (30 Jul 2023)-img/p2_chem_s4_q08_4.png}{164.3}{0}
\qfigdef{p2_chem_s4_q08_5}{TAAS - Test 02 (30 Jul 2023)-img/p2_chem_s4_q08_5.png}{127.3}{0}
\raggedcolumns
\makeatletter\@beginparpenalty=10000\makeatother
\newcolumntype{L}[1]{>{\raggedright\arraybackslash}p{#1}}
\begin{document}
\begin{center}{\Large\bfseries TAAS --- Test 02 (30 Jul 2023)}\end{center}
\begin{multicols}{2}
\raggedcolumns
\papersec{1}
\partsec{1}{PHYSICS}
\sect{SECTION-I(i)}
\qn{1} For an infinite line of charge having charge density $\lambda$ lying along $x$-axis, the work required in moving charge $q$ from $C$ to $A$ along arc $CA$ is :-
\qfig{p1_phy_s1_q01_1}
\begin{qoptions}
\opt{A}{$\dfrac{q\lambda}{\pi\varepsilon_0}\log_e\sqrt2$}
\opt{B}{$\dfrac{q\lambda}{4\pi\varepsilon_0}\log_e\sqrt2$}
\opt{C}{$\dfrac{q\lambda}{4\pi\varepsilon_0}\log_e2$}
\opt{D}{$\dfrac{q\lambda}{2\pi\varepsilon_0}\log_e\dfrac12$}
\end{qoptions}
\qn{2} In the circuit diagram shown all the capacitors are in $\mu$F. The equivalent capacitance (in $\mu$F) between points $A$ and $B$ is :-
\qfig{p1_phy_s1_q02_1}
\begin{qoptions}
\opt{A}{$\dfrac{14}5$}
\opt{B}{$\dfrac75$}
\opt{C}{$\dfrac37$}
\opt{D}{None of these}
\end{qoptions}
\qn{3} Ratio of charges on 5$\mu$F and 10$\mu$F capacitor when circuit is in steady state :-
\qfig{p1_phy_s1_q03_1}
\begin{qoptions}
\opt{A}{38 : 13}
\opt{B}{19 : 26}
\opt{C}{13 : 38}
\opt{D}{26 : 19}
\end{qoptions}
\qn{4} A particle of positive charge $q$ and mass $m$ enters with velocity $V\hat j$ at the origin in a magnetic field $B\left(-\hat k\right)$ which is present in the whole space. The charge makes a perfectly inelastic collision with free identical particle at rest at its maximum $y$-coordinate. After collision the combined charge will move on trajectory: (where $r=\dfrac{mV}{qB}$) (Neglect earth's gravity)
\begin{qoptions}
\opt{A}{$y=\dfrac{mV}{qB}\left(-\hat i\right)$}
\opt{B}{$(x+r)^2+\left(y-\dfrac r2\right)^2=\dfrac{r^2}4$}
\opt{C}{$(x-r)^2+(y-r)^2=r^2$}
\opt{D}{$(x-r)^2+\left(y+\dfrac r2\right)^2=\dfrac{r^2}4$}
\end{qoptions}
\sect{SECTION-I(ii)}
\qn{5} Two gases have same initial pressure, volume, temperature. Both gases are compressed to same volume one isothermally and one adiabatically :-
\begin{qoptions}
\opt{A}{The final temperature is greater for the adiabatic process.}
\opt{B}{The final pressure is greater for the adiabatic process.}
\opt{C}{The work done on the gas is greater for the isothermal process.}
\opt{D}{All the above options are incorrect.}
\end{qoptions}
\qn{6} A cyclic process $ABCD$ is shown in the $P$-$V$ diagram. Which of the following curves represents the same process if $BC$ and $DA$ are isothermal processes.
\qfig{p1_phy_s2_q06_1}
\begin{qoptions}
\opt{A}{\qfigi{p1_phy_s2_q06_A}}
\opt{B}{\qfigi{p1_phy_s2_q06_B}}
\opt{C}{\qfigi{p1_phy_s2_q06_C}}
\opt{D}{\qfigi{p1_phy_s2_q06_D}}
\end{qoptions}
\qn{7} A conducting rod of length $\ell$ moves with velocity $u$ in a direction parallel to a long wire carrying a steady current $I$. The axis of the rod is maintained perpendicular to the wire with near end a distance $r$ away as shown in the fig. Find the e.m.f. induced in the rod.
\qfig{p1_phy_s2_q07_1}
\begin{qoptions}
\opt{A}{$-\dfrac{\mu_0Iu}{2\pi}\ell n\left(\dfrac r{r+\ell}\right)$}
\opt{B}{$\dfrac{2\mu_0Iu}\pi\ell n\left(\dfrac{r+\ell}r\right)$}
\opt{C}{$\dfrac{-\mu_0Iu}\pi\ell n\left(\dfrac\ell{r+\ell}\right)$}
\opt{D}{$\dfrac{\mu_0Iu}{2\pi}\ell n\left(\dfrac{r+\ell}r\right)$}
\end{qoptions}
\sect{SECTION-I(iii)}
\qn{8} A particle is taken to a distance $r$ ($>R$) from centre of earth. $R$ is the radius of the earth. It is given velocity $v$ which is perpendicular to radius. With the given values of $v$ in List-I, you have to match the resultant path of particle in List-II.
\begin{center}\footnotesize\setlength{\tabcolsep}{2pt}
\begin{tabular}{|L{0.6cm}L{2.5cm}|L{0.6cm}L{1.8cm}|}\hline
\multicolumn{2}{|c|}{List-I}&\multicolumn{2}{c|}{List-II}\\\hline
(P)&$v=\sqrt{\dfrac{GM}r}$&(1)&Elliptical\\\hline
(Q)&$v=\sqrt{\dfrac{2GM}r}$&(2)&Parabolic\\\hline
(R)&$v>\sqrt{\dfrac{2GM}r}$&(3)&Hyperbolic\\\hline
(S)&$\sqrt{\dfrac{GM}r}<v<\sqrt{\dfrac{2GM}r}$&(4)&Circular\\\hline
&&(5)&Straight line\\\hline
\end{tabular}\end{center}
\begin{qoptions}
\opt{A}{P $\to$ 3;Q $\to$ 1;R $\to$ 4;S $\to$ 2}
\opt{B}{P $\to$ 2;Q $\to$ 4;R $\to$ 3;S $\to$ 1}
\opt{C}{P $\to$ 4;Q $\to$ 2;R $\to$ 3;S $\to$ 1}
\opt{D}{P $\to$ 1;Q $\to$ 2;R $\to$ 3;S $\to$ 4}
\end{qoptions}
\qn{9} Magnetic flux in a circular coil of resistance 100 $\Omega$ changes with time as shown in figure. Then match the following list (inside magnetic flux is positive) :-
\qfig{p1_phy_s3_q09_1}
\begin{center}\footnotesize\setlength{\tabcolsep}{2pt}
\begin{tabular}{|L{0.6cm}L{2.5cm}|L{0.6cm}L{2.5cm}|}\hline
\multicolumn{2}{|c|}{List-I}&\multicolumn{2}{c|}{List-II}\\\hline
(P)&At 1 second induced current is&(1)&Clockwise\\\hline
(Q)&At 5 second induced current is&(2)&Anticlockwise\\\hline
(R)&At 9 second induced current is&(3)&Zero\\\hline
(S)&At 15 second induced current is&(4)&Magnitude is linearly increasing\\\hline
&&(5)&Magnitude is linearly decreasing\\\hline
\end{tabular}\end{center}
\begin{qoptions}
\opt{A}{P $\to$ 1,3;Q $\to$ 2,4;R $\to$ 1,5;S $\to$ 2}
\opt{B}{P $\to$ 2,5;Q $\to$ 3;R $\to$ 1,5;S $\to$ 2}
\opt{C}{P $\to$ 1,5;Q $\to$ 4;R $\to$ 2,5;S $\to$ 3}
\opt{D}{P $\to$ 2,5;Q $\to$ 2,4;R $\to$ 3,5;S $\to$ 1,2}
\end{qoptions}
\qn{10}
\qfig{p1_phy_s3_q10_1}
\begin{center}\footnotesize\setlength{\tabcolsep}{2pt}
\begin{tabular}{|L{0.6cm}L{3.3cm}|L{0.6cm}L{0.8cm}|}\hline
\multicolumn{2}{|c|}{List-I}&\multicolumn{2}{c|}{List-II}\\\hline
(P)&Current through battery just after closing of switch (in ampere)&(1)&12\\\hline
(Q)&Current through capacitor just after closing of switch (in ampere)&(2)&8\\\hline
(R)&Switch is closed and when steady state of current is achieved the current through battery (in ampere)&(3)&9\\\hline
(S)&After steady state switch is re-opened just after reopening of switch the ratio of current through inductor and capacitor&(4)&1\\\hline
&&(5)&0\\\hline
\end{tabular}\end{center}
\begin{qoptions}
\opt{A}{P $\to$ 1;Q $\to$ 2;R $\to$ 3;S $\to$ 4}
\opt{B}{P $\to$ 3;Q $\to$ 4;R $\to$ 2;S $\to$ 1}
\opt{C}{P $\to$ 4;Q $\to$ 3;R $\to$ 1;S $\to$ 2}
\opt{D}{P $\to$ 2;Q $\to$ 1;R $\to$ 4;S $\to$ 3}
\end{qoptions}
\qn{11} In all the situations current in loop-I is $i_1$ and current in loop-II is $i_2$. Consider the infinite wire as the side of infinite large loop. Column-II describes the variation in current $i_1$ in different arrangements and column-I describes the various effects.
\begin{center}\footnotesize\setlength{\tabcolsep}{2pt}
\begin{tabular}{|L{0.6cm}L{2.0cm}|L{0.6cm}L{3.8cm}|}\hline
\multicolumn{2}{|c|}{List-I}&\multicolumn{2}{c|}{List-II}\\\hline
(P)&Current is clockwise in loop-II.&(1)&\qfigi{p1_phy_s3_q11_1}\par Current $i_1$ is clockwise and decreasing at constant rate.\\\hline
(Q)&Current is anticlockwise in loop-II.&(2)&\qfigi{p1_phy_s3_q11_2}\par Current $i_1$ is clockwise and decreasing at constant rate.\\\hline
(R)&Flux of current $i_1$ through loop-II is less than flux of current $i_1$ through loop-I&(3)&\qfigi{p1_phy_s3_q11_3}\par Current in infinite wire is $i_1$ and decreasing at constant rate. The wire is perpendicular to plane of loop.\\\hline
(S)&Loop-II tends to reduce its area due to magnetic force applied by magnetic field of $i_1$&(4)&\qfigi{p1_phy_s3_q11_4}\par Loop-I having constant current in clockwise direction moving upward with retardation. Both loops are co-axial.\\\hline
&&(5)&\qfigi{p1_phy_s3_q11_5}\par Current in infinite wire is $i_1$ upward and decreasing at constant rate. The wire is parallel to the plane of loop.\\\hline
\end{tabular}\end{center}
\begin{qoptions}
\opt{A}{P $\to$ 2,5;Q $\to$ 1,4;R $\to$ 1,2,4,5;S $\to$ 2,4}
\opt{B}{P $\to$ 2,5;Q $\to$ 1,3,4;R $\to$ 1,3,4,5;S $\to$ 2,4}
\opt{C}{P $\to$ 2,4;Q $\to$ 1,3,4;R $\to$ 1,2,3,4,5;S $\to$ 2,5}
\opt{D}{P $\to$ 2,4;Q $\to$ 1,4;R $\to$ 1,2,3,4,5;S $\to$ 2,5}
\end{qoptions}
\sect{SECTION-II}
\qn{1} $n$-moles of an ideal gas at $27^\circ$C are expanded isothermally to five time its volume and then heated at constant volume to raise its pressure equal to initial pressure (before expansion). If heat given to the system in the process is 83.14 kJ then find the value of $n$ (to the nearest integer). The $\gamma\left(=\dfrac{C_P}{C_V}\right)$ of the gas is 1.42. [$\ell$n5 $=1.61$]
\qn{2} Two capacitors A and B are connected in series across a 100V supply and it is observed that the potential difference across them are 60V and 40V. A capacitor of 2$\mu$F capacitance is now connected in parallel with A and the potential difference across B raises to 90 V. If capacitance of A is 0.04 x$\mu$F. Find $x$.
\qn{3} A resistance coil wired to an external battery is placed inside a thermally insulated cylinder fitted with a frictionless piston and containing an ideal gas. A current $I=2.40$ mA flows through the coil, which has a resistance $R=490\ \Omega$. If the speed $v$ with which piston of mass $m=12$ kg move upward, so that the temperature of the gas remains unchanged is 0.003$x$ mm/sec. Find $x$?
\qfig{p1_phy_s4_q03_1}
\qn{4} The circular current loop of radius $b=2$cm shown in the figure is mounted rigidly on the axle, midway between the two supporting cords. Current in the loop is 100 amp. The length of axle is $\ell=1$ m. Absence of an external magnetic field the tensions in the cords are equal to $T_0=10\pi$N. When the vertical magnetic field $B=100$ T is present and the axle still in equilibrium, then the tensions (in newtons) in the cords are $T_1=n\pi$; $T_2=m\pi$. Find the magnitude of $m-n$.
\qfig{p1_phy_s4_q04_1}
\qn{5} Two parallel long smooth conducting rails separated by a distance $\ell=10$ cm are connected by a movable conducting connector of mass $m=4$ mg. Terminals of the rails are connected by the resistor $R=2\Omega$ and capacitance $C=1\mu$F as shown. A uniform magnetic field $B=20$ T perpendicular to the plane of the rails is switched on. The connector is dragged by a constant force $F=10$ N. The speed of the connector as function of time if the force $F$ is applied at $t=0$ is equal to $v=5(1-e^{-x\times10^4\times t})$ m/s. Find the value of $x$.
\qfig{p1_phy_s4_q05_1}
\qn{6} A cavity of radius $r$ is present inside a fixed solid dielectric sphere of radius $R$, having a volume charge density of $\rho$. The distance between the centres of the sphere and the cavity is $a$. An electron $e$ is kept inside the cavity at an angle $\theta=45^\circ$ as shown. The electron (of mass $m$ and charge $-e$) will take $\left(\dfrac{P\sqrt2mr\varepsilon_0}{ea\rho}\right)^{1/2}$ time to touch the sphere again. Neglect gravity. Find the value of $P$.
\qfig{p1_phy_s4_q06_1}
\partsec{2}{CHEMISTRY}
\sect{SECTION-I(i)}
\qn{1} (I). Salt A $+$ dil. H$_2$SO$_4\rightarrow$ ``B'' gas [B has suffocating order]

(II). Salt A $+$ Pb(OAc)$_2\rightarrow$ C($\downarrow$) $\xrightarrow{\Delta/\mathrm{O_2}}$ D($\downarrow$) [C \& D both are white ppt.]

Have A and D are :-
\begin{qoptions}
\opt{A}{HSO$_3^-$ , PbSO$_4$}
\opt{B}{SO$_3^{2-}$ , PbSO$_3$}
\opt{C}{SO$_3^{2-}$ , PbSO$_4$}
\opt{D}{CO$_3^{2-}$ , PbSO$_3$}
\end{qoptions}
\qn{2} \qfigi{p1_chem_s1_q02_1}$\xrightarrow{\mathrm{N.B.S.}}$ (P) $\xrightarrow[\text{ether}]{\mathrm{Mg}}$ (Q) $\xrightarrow[\mathrm{H_2O/H^+}]{\mathrm{CO_2}}$ (R) $\xrightarrow{\mathrm{L.A.H.}}$ (S) $\xrightarrow{\mathrm{PCC}}$ (T)

Which of the following is correct ?
\begin{qoptions}
\opt{A}{D.U. of (R) is 2}
\opt{B}{(S) is 2$^\circ$ alcohol}
\opt{C}{(T) have four sp$^3$ hybrid carbon.}
\opt{D}{Reaction of (R) with sodalime gives methyl cyclohexene.}
\end{qoptions}
\qn{3} \qfig{p1_chem_s1_q03_1}
Correct statement is?
\begin{qoptions}
\opt{A}{(C) on reaction with NH$_2$OH give compound which show G.I.}
\opt{B}{Number of $\pi$-bond present in (C) are 7.}
\opt{C}{(C) is an aldehyde}
\opt{D}{Compound (C) show tautomerism.}
\end{qoptions}
\qn{4} 20 mL of H$_2$O$_2$ is reacted completely with acidified K$_2$Cr$_2$O$_7$ solution. 40 mL of K$_2$Cr$_2$O$_7$ was required to oxidise the H$_2$O$_2$ completely. Also, 2.0 mL of the same K$_2$Cr$_2$O$_7$ solution required 5.0 mL of a 1.0 M H$_2$C$_2$O$_4$ solution to reach equivalence point. Which of the following statements is/are correct ?
\begin{qoptions}
\opt{A}{The H$_2$O$_2$ solution is 5M.}
\opt{B}{The volume strength of H$_2$O$_2$ is 28 V.}
\opt{C}{The volume strength of H$_2$O$_2$ is 9.46 V.}
\opt{D}{If 40 mL of 5/8 M H$_2$O$_2$ is further added to the 10 mL of above H$_2$O$_2$ solution the volume strength of the resulting solution is changed to 1.5 V.}
\end{qoptions}
\sect{SECTION-I(ii)}
\qn{5} Consider the following parallel, first order processes
\qfigi{p1_chem_s2_q05_1} $k_1$, $k_2$ are rate constants and $E_1$ and $E_2$ are activation energies.

The correct deduction(s) is/are :
\begin{qoptions}
\opt{A}{At any time, the ratio of concentration of P and Q has a fixed value at a given temperature.}
\opt{B}{If a catalyst is added, it will always decrease both $E_1$ and $E_2$ but by same amount.}
\opt{C}{If `$E$' represents the overall activation energy for decay of A, it will be $\left(\dfrac{E_1k_1+E_2k_2}{k_1+k_2}\right)$.}
\opt{D}{If `$k$' represents the overall rate constant for decay of A, it will be $\dfrac1k=\dfrac1{k_1}+\dfrac1{k_2}$.}
\end{qoptions}
\qn{6} The vapour pressure of solid benzoic acid ($P$ in bar) has been found to obey the relationship (in the neighbourhood of 298 K) as : $\ln\dfrac P{P^0}=a-\dfrac bT$ ; where $a=20$ and $b=1\times10^4$ K. Which of the following values are correct at 298 K for the sublimation of benzoic acid : ($R=8.3$ J/mole-K) ($P^\circ=1$ bar)
\begin{qoptions}
\opt{A}{$\Delta G^\circ=33.53$ kJ mol$^{-1}$}
\opt{B}{$\Delta H^\circ=83$ kJ mol$^{-1}$}
\opt{C}{$\Delta S^\circ=166$ JK$^{-1}$ mol$^{-1}$}
\opt{D}{$\Delta H^\circ=-88.96$ kJ mol$^{-1}$}
\end{qoptions}
\qn{7} Which is/are correct statement -
\begin{qoptions}
\opt{A}{$\mathrm{H_2(g)+\tfrac12O_2(g)\to H_2O(l)}$ ; $\Delta_r H$

$\Delta_rH=\Delta H_C(\mathrm{H_2,g})=\Delta H_f(\mathrm{H_2O},\ell)$}
\opt{B}{Na(s)$\to$ Na(g) ; $\Delta_rH$

$\Delta_rH=\Delta H_{\text{sublimation}}(\mathrm{Na,s})=\Delta H_{\text{atomisation}}(\mathrm{Na,s})$}
\opt{C}{C(graphite) $+$ O$_2$(g)$\to$ CO$_2$(g) ; $\Delta_rH$

$\Delta_rH=\Delta H_C(\mathrm{C,graphite})=\Delta H_f(\mathrm{CO_2,g})$}
\opt{D}{H$_2$(g)$\to$ 2H(g) ; $\Delta_rH$

$\Delta_rH=\Delta H_{\text{atomisation}}(\mathrm{H_2,g})=\Delta H_f(\mathrm{H,g})$}
\end{qoptions}
\sect{SECTION-I(iii)}
\qn{8} Match the List-I with List-II
\begin{center}\footnotesize\setlength{\tabcolsep}{2pt}
\begin{tabular}{|L{0.6cm}L{1.5cm}L{1.5cm}|L{0.6cm}L{1.0cm}|}\hline
\multicolumn{3}{|L{3.6cm}|}{\centering List-I\\(Colour of borax bead test (cold)}&\multicolumn{2}{L{1.6cm}|}{\centering List-II\\(Metal present)}\\\hline
&Oxidising flame&Reducing flame&&\\\hline
(P)&Blue&Blue&(1)&Cu\\\hline
(Q)&Violet&Colorless&(2)&Fe\\\hline
(R)&Yellow&Green&(3)&Mn\\\hline
(S)&Blue&Red Opaqe&(4)&Co\\\hline
&&&(5)&Cr\\\hline
\end{tabular}\end{center}
The correct option is :
\begin{qoptions}
\opt{A}{P $\to$ 3;Q $\to$ 4;R $\to$ 1;S $\to$ 3}
\opt{B}{P $\to$ 4;Q $\to$ 3;R $\to$ 1;S $\to$ 2}
\opt{C}{P $\to$ 4;Q $\to$ 2;R $\to$ 3;S $\to$ 1}
\opt{D}{P $\to$ 4;Q $\to$ 3;R $\to$ 2;S $\to$ 1}
\end{qoptions}
\qn{9} Consider the following extraction process and match List-I with List-II.
\qfig{p1_chem_s3_q09_1}
\begin{center}\footnotesize\setlength{\tabcolsep}{2pt}
\begin{tabular}{|L{0.6cm}L{2.5cm}|L{0.6cm}L{1.8cm}|}\hline
\multicolumn{2}{|c|}{List-I}&\multicolumn{2}{c|}{List-II}\\\hline
(P)&Magnetic separation&(1)&Step-I\\\hline
(Q)&Liquation&(2)&Step-II\\\hline
(R)&Poling&(3)&Step-III\\\hline
(S)&Carbon reduction&(4)&Step-IV\\\hline
&&(5)&No any step\\\hline
\end{tabular}\end{center}
The correct option is :
\begin{qoptions}
\opt{A}{P $\to$ 4;Q $\to$ 1;R $\to$ 2;S $\to$ 3}
\opt{B}{P $\to$ 4;Q $\to$ 1;R $\to$ 3;S $\to$ 2}
\opt{C}{P $\to$ 1;Q $\to$ 4;R $\to$ 3;S $\to$ 2}
\opt{D}{P $\to$ 1;Q $\to$ 4;R $\to$ 2;S $\to$ 3}
\end{qoptions}
\qn{10}
\begin{center}\footnotesize\setlength{\tabcolsep}{2pt}
\begin{tabular}{|L{0.6cm}L{3.3cm}|L{0.6cm}L{2.2cm}|}\hline
\multicolumn{2}{|L{3.9cm}|}{\centering List-I\\(Reaction)}&\multicolumn{2}{L{2.8cm}|}{\centering List-II\\(Major product)}\\\hline
(P)&\qfigi{p1_chem_s3_q10_1}&(1)&1$^\circ$ alcohol\\\hline
(Q)&\qfigi{p1_chem_s3_q10_2}&(2)&2$^\circ$ alcohol\\\hline
(R)&\qfigi{p1_chem_s3_q10_3}&(3)&3$^\circ$ alcohol\\\hline
(S)&CH$_3{-}$C$\equiv$N $+$ CH$_3$MgBr $\to$&(4)&Ketone\\\hline
&&(5)&Product show O.I.\\\hline
\end{tabular}\end{center}
Match the correct code
\begin{qoptions}
\opt{A}{P $\to$ 2,5;Q $\to$ 3;R $\to$ 1;S $\to$ 4}
\opt{B}{P $\to$ 2,5;Q $\to$ 3,5;R $\to$ 1;S $\to$ 4}
\opt{C}{P $\to$ 2;Q $\to$ 3,5;R $\to$ 1;S $\to$ 4}
\opt{D}{P $\to$ 1;Q $\to$ 3;R $\to$ 1;S $\to$ 3}
\end{qoptions}
\qn{11} One mole of a monatomic ideal gas is taken along two cyclic processes $E\to F\to G\to E$ and $E\to F\to H\to E$ as shown in the $p$-$V$ diagram.
\qfig{p1_chem_s3_q11_1}
The processes involved are purely isochoric, isobaric, isothermal ($F\to G$) or adiabatic ($F\to H$). Match the paths in List-I with magnitudes of the work done in List-II and select the correct answer using the codes given below the lists.
\begin{center}\footnotesize\setlength{\tabcolsep}{2pt}
\begin{tabular}{|L{0.6cm}L{1.1cm}|L{0.6cm}L{2.4cm}|}\hline
\multicolumn{2}{|c|}{List-I}&\multicolumn{2}{c|}{List-II}\\\hline
(P)&$G\to E$&(1)&$160\,p_0V_0\ln2$\\\hline
(Q)&$G\to H$&(2)&$36\,p_0V_0$\\\hline
(R)&$F\to H$&(3)&$24\,p_0V_0$\\\hline
(S)&$F\to G$&(4)&$31\,p_0V_0$\\\hline
&&(5)&$96\,p_0V_0$\\\hline
\end{tabular}\end{center}
\begin{qoptions}
\opt{A}{P $\to$ 4;Q $\to$ 3;R $\to$ 5;S $\to$ 2}
\opt{B}{P $\to$ 3;Q $\to$ 2;R $\to$ 4;S $\to$ 1}
\opt{C}{P $\to$ 4;Q $\to$ 3;R $\to$ 2;S $\to$ 1}
\opt{D}{P $\to$ 1;Q $\to$ 2;R $\to$ 3,5;S $\to$ 4}
\end{qoptions}
\sect{SECTION-II}
\qn{1} 
\par\noindent\qfigi{p1_chem_s4_q01_3}$\xrightarrow[\hphantom{xx}]{\mathrm{mCPBA}}\xrightarrow{\mathrm{H_2O/H^+}}$ ($x$) number of organic product (Vicinal diol)
\par\noindent\qfigi{p1_chem_s4_q01_4}$\xrightarrow[\text{Reagent}]{\text{Bayers}}$ ($y$) number of organic product (Vicinal diol)
\par\noindent\qfigi{p1_chem_s4_q01_1}$\xrightarrow{\mathrm{RCO_3H/H_2O}}$ ($z$) number of organic product (Vicinal diol)
\par\noindent\qfigi{p1_chem_s4_q01_2}$\xrightarrow{\mathrm{OsO_4/NaHSO_3}}$ ($k$) number of organic product (Vicinal diol)
\par What is the value of $x+y+z+k$?
\qn{2} How many different organic products are formed by Ozonolysis of ortho xylene?
\qn{3} \qfig{p1_chem_s4_q03_1}
Calculate number of lone pair in product (z) ?
\qn{4} Find number of reaction in which at least one of the atom is in zero oxidation state in product side.

(i) $\mathrm{MnO_4^-/H^++H_2O_2\rightarrow}$

(ii) 2KI$_\text{(aq.)}$ $+$ Cl$_2$ water $\rightarrow$

(iii) OCl$^-$ $+$ Cl$^-$ $\xrightarrow[\text{solution}]{\text{acidic}}$

(iv) KI$_3$ $+$ S$_2$O$_3^{2-}\rightarrow$

(v) SO$_{2\text{(aq.)}}$ $+$ H$_2$S $\rightarrow$

(vi) Fe$^{3+}_{\text{(aq.)}}$ $+$ H$_2$S $\xrightarrow[\text{solution}]{\text{acidic}}$

(vii) NH$_4$NO$_3$ $\xrightarrow[250^\circ\mathrm{C}]{\Delta}$
\qn{5} \qfig{p1_chem_s4_q05_1}
Find sum of metal$-$oxygen linkages in X, Y and Z.
\qn{6} An element X (Atomic weight $=24$ gm/mole) form FCC lattice. If edge length of unit cell is $4\times10^{-8}$ cm and observed density is $2.4\times10^3$ kg/m$^3$, then percentage occupancy of lattice points of X element is (Given $N_A=6\times10^{23}$)
\columnbreak
\partsec{3}{MATHEMATICS}
\sect{SECTION-I(i)}
\qn{1} Let $f(x)=x^2+xg'(1)+g''(2)$ and $g(x)=f(1)x^2+xf'(x)+f''(x)$, then the value of $g'(1)+g''(2)$ is -
\begin{qoptions}
\opt{A}{3}
\opt{B}{$-3$}
\opt{C}{0}
\opt{D}{2}
\end{qoptions}
\qn{2} $\displaystyle\int\dfrac{\sin2x+\sin4x-\sin6x}{1+\cos2x+\cos4x+\cos6x}dx$ is equal to (where $C$ is constant of integration)
\begin{qoptions}
\opt{A}{$\dfrac13\ell n|\sec3x|+\dfrac12\ell n|\sec2x|+\ell n|\sec x|+C$}
\opt{B}{$\dfrac13\ell n|\sec3x|-\dfrac12\ell n|\sec2x|-\ell n|\sec x|+C$}
\opt{C}{$\dfrac13\ell n|\sec3x|-\dfrac12\ell n|\sec2x|+\ell n|\sec x|+C$}
\opt{D}{$\dfrac13\ell n|\sec3x|+\dfrac12\ell n|\sec2x|-\ell n|\sec x|+C$}
\end{qoptions}
\qn{3} $f(x)=\tan(\sin^{-1}(\cos(\sin^{-1}(x))))$ $\tan(\cos^{-1}(\sin(\cos^{-1}x)))$ where $x\in(0,1)$ then $\displaystyle\sum_{r=0}^{\infty}\dfrac1{2^r}f\left(\dfrac1{2^r}\right)$ is equal to
\begin{qoptions}
\opt{A}{$\dfrac12$}
\opt{B}{$\dfrac52$}
\opt{C}{3}
\opt{D}{2}
\end{qoptions}
\qn{4} Let a real valued differentiable function $f(x)$ be strictly monotonic on $[a,b]$, then
\begin{qoptions}
\opt{A}{$\displaystyle\int_a^b\left[(f(x))^2-(f(b))^2\right]dx=$\par\noindent$\displaystyle\int_{f(a)}^{f(b)}2x\left(b-f^{-1}(x)\right)dx$}
\opt{B}{$\displaystyle\int_a^b\left[(f(x))^2-(f(a))^2\right]dx=$\par\noindent$\displaystyle\int_{f(a)}^{f(b)}2x\left(b-f^{-1}(x)\right)dx$}
\opt{C}{$\displaystyle\int_a^b[f(x)-f(a)]^2dx=\int_{f(a)}^{f(b)}x\left(b-f^{-1}(x)\right)^2dx$}
\opt{D}{None of these}
\end{qoptions}
\sect{SECTION-I(ii)}
\qn{5} Let $f(x)=\displaystyle\int_x^{x+\frac\pi3}|\sin\theta|\,d\theta$; where $x\in[0,\pi]$, which of the following is/are true
\begin{qoptions}
\opt{A}{maximum value of $f(x)$ is 1}
\opt{B}{minimum value of $f(x)$ is $\dfrac32$}
\opt{C}{maximum value of $f(x)$ is $\dfrac12$}
\opt{D}{minimum value of $f(x)$ is $2-\sqrt3$}
\end{qoptions}
\qn{6} Let $f:(a,b)\to R$ is a differentiable function such that $\displaystyle\lim_{x\to a^+}f^2(x)=0$, $\displaystyle\lim_{x\to b^-}f^2(x)=e-1$ and $2f(x)\,f'(x)-f^2(x)\geq1$ for all $x\in(a,b)$ then value of $(b-a)$ can be
\begin{qoptions}
\opt{A}{0}
\opt{B}{$\dfrac12$}
\opt{C}{1}
\opt{D}{2}
\end{qoptions}
\qn{7} Let $f(x)=\cos^{-1}\left(\ell n\left(\dfrac{x+1}x\right)\right)$, then which of the following statement(s) is/are correct
\begin{qoptions}
\opt{A}{Domain of $f(x)$ contains $(-\infty,-1)\cup(1,\infty)$}
\opt{B}{Range of $f(x)$ contains $\left[0,\dfrac1{e-1}\right]$}
\opt{C}{Domain of $f(x)$ is $\left(-\infty,\dfrac e{1-e}\right]\cup\left[\dfrac1{e-1},\infty\right)$}
\opt{D}{Range of $f(x)$ is $(-\infty,-1]\cup(1,\infty)$}
\end{qoptions}
\sect{SECTION-I(iii)}
\qn{8} Match List-I with List-II and select the correct answer using the code given below the list.
\begin{center}\footnotesize\setlength{\tabcolsep}{2pt}
\begin{tabular}{|L{0.6cm}L{4.5cm}|L{0.8cm}|}\hline
\multicolumn{2}{|c|}{List-I}&List-II\\\hline
(P)&If $x=\alpha$ satisfies $\sin^{-1}x+\cos^{-1}x^2+\dfrac\pi2=0$, then the value of $\dfrac{\sec^{-1}\alpha-\tan^{-1}\alpha}{\cot^{-1}\alpha-\cos\mathrm{ec}^{-1}\alpha}$ is&(1) 0\\\hline
(Q)&If $a=\displaystyle\lim_{x\to\infty}(2.3^x+3.4^x+4.5^x)^{\frac1x}$ and $b=\displaystyle\lim_{x\to(-\infty)}(2.3^x+3.4^x+4.5^x)^{\frac1x}$ then $\left(\dfrac{a+b}{a-b}\right)$ is equal to&(2) 4\\\hline
(R)&Number of points of discontinuity of $f(x)=x\sin(\pi(2x-[x]))$ in $[0,2]$ will be (Note $[k]$ denotes greatest integer function less than or equal to $k$)&(3) 1\\\hline
(S)&$\displaystyle\lim_{x\to0}\left(\dfrac{3e^{x/3}-x-3}{x^2}\right)$ is equal to&(4) $\dfrac13$\\\hline
&&(5) $\dfrac16$\\\hline
\end{tabular}\end{center}
\begin{qoptions}
\opt{A}{P $\to$ 1;Q $\to$ 2;R $\to$ 2;S $\to$ 4}
\opt{B}{P $\to$ 3;Q $\to$ 1;R $\to$ 2;S $\to$ 5}
\opt{C}{P $\to$ 3;Q $\to$ 2;R $\to$ 1;S $\to$ 5}
\opt{D}{P $\to$ 1;Q $\to$ 3;R $\to$ 2;S $\to$ 4}
\end{qoptions}
\qn{9} For a continuous function $f(x)$, equation $f(x)=0$ has a 4 distinct solutions then match List-I with List-II and select the correct answer using the code given below the list.
\begin{center}\footnotesize\setlength{\tabcolsep}{2pt}
\begin{tabular}{|L{0.6cm}L{1.8cm}|L{0.7cm}L{3.5cm}|}\hline
\multicolumn{2}{|c|}{List-I}&\multicolumn{2}{c|}{List-II}\\\hline
(P)&Maximum number of distinct solutions of $f(|x|)=0$ is equal to the value of&(1)&$\displaystyle\lim_{x\to0^+}\left(\dfrac{x+1}{2x^2+3x+5}\right)^{1/x}$\\\hline
(Q)&Minimum number of distinct solutions of $f(|x|)=0$ is equal to the value of&(2)&$\displaystyle\lim_{x\to0}\left[\dfrac{9\sin x}x\right]$ (where $[.]$ is greatest integer function)\\\hline
(R)&Maximum number of distinct solution of $f(-|x|)=0$ is equal to the value of&(3)&$\displaystyle\int_2^{18}\dfrac{f(x)}{f(20-x)+f(x)}dx$\\\hline
(S)&Minimum number of distinct solution of $f(-|x|)=0$ is equal to the value of&(4)&$\displaystyle\int_{-2}^2\dfrac{x^3-x}{\sqrt{4-x^2}}dx$\\\hline
&&(5)&$24\displaystyle\int_{1/3}^{2/3}\dfrac{\ell nx}{\ell n(x-x^2)}dx$\\\hline
\end{tabular}\end{center}
\begin{qoptions}
\opt{A}{P $\to$ 3,4;Q $\to$ 1,2,4;R $\to$ 2,3;S $\to$ 1,2,4}
\opt{B}{P $\to$ 2,5;Q $\to$ 1,4,5;R $\to$ 2,3;S $\to$ 1,2,4}
\opt{C}{P $\to$ 2,3;Q $\to$ 1,4;R $\to$ 2,3;S $\to$ 1,4}
\opt{D}{P $\to$ 1,2,3;Q $\to$ 1,3;R $\to$ 1,2,3;S $\to$ 1,2,4}
\end{qoptions}
\qn{10} Match List-I with List-II and select the correct answer using the code given below the list.
\begin{center}\footnotesize\setlength{\tabcolsep}{2pt}
\begin{tabular}{|L{0.6cm}L{3.9cm}|L{0.7cm}L{0.6cm}|}\hline
\multicolumn{2}{|c|}{List-I}&\multicolumn{2}{c|}{List-II}\\\hline
(P)&Value of $\dfrac8{\pi^2}\displaystyle\int_0^\pi\sin^{-1}(\sin x)\,dx$ is equal to&(1)&0\\\hline
(Q)&Value of $\displaystyle\int_0^\pi\cos mx\,\sin nx\,dx$ is equal to (Here $m$ and $n\in I$ \& $(m-n)$ is a non-zero even integer)&(2)&1\\\hline
(R)&Value of $\displaystyle\lim_{x\to0^+}\sin\cos^{-1}\cot^{-1}\left(\dfrac1x\right)$ is equal to&(3)&2\\\hline
(S)&$\displaystyle\lim_{x\to2^+}\left(\dfrac{2\{x\}-4}{[x]-3}\right)$ is equal to (Here $[.]$ \& $\{.\}$ denote greatest integer \& fractional part function)&(4)&4\\\hline
&&(5)&3\\\hline
\end{tabular}\end{center}
\begin{qoptions}
\opt{A}{P $\to$ 3;Q $\to$ 2;R $\to$ 1;S $\to$ 4}
\opt{B}{P $\to$ 2;Q $\to$ 3;R $\to$ 5;S $\to$ 1}
\opt{C}{P $\to$ 3;Q $\to$ 1;R $\to$ 2;S $\to$ 4}
\opt{D}{P $\to$ 5;Q $\to$ 2;R $\to$ 4;S $\to$ 1}
\end{qoptions}
\qn{11} Match the statement in List-I with given value in List-II
\begin{center}\footnotesize\setlength{\tabcolsep}{2pt}
\begin{tabular}{|L{0.6cm}L{4.4cm}|L{0.7cm}L{0.4cm}|}\hline
\multicolumn{2}{|c|}{List-I}&\multicolumn{2}{c|}{List-II}\\\hline
(P)&If area bounded by $x=(\theta-\sin\theta)$; $y=(1-\cos\theta)$; $0\leq\theta\leq2\pi$ is $\lambda\pi$; then the value of $\lambda$ can be equal to&(1)&6\\\hline
(Q)&If the least area of triangle formed by tangent to the circle $x^2+y^2=1$ and $x=0$, $y=0$ is $A$, then $A$ is coprime with&(2)&3\\\hline
(R)&Area of triangle formed by any tangent to the curve $y=\dfrac{2x^2-x-1}{2x^2+5x+2}$ $\left(x\neq-2,-\dfrac12\right)$ and pair of lines $xy-x+2y=2$ is&(3)&80\\\hline
(S)&If the maximum area bounded by the curves $x^2=4ay$, $y=ax$, $y=\dfrac xa$, $1\leq a\leq2$ is $\alpha$, then value of $(\alpha+a)$ is&(4)&84\\\hline
&&(5)&86\\\hline
\end{tabular}\end{center}
\begin{qoptions}
\opt{A}{P $\to$ 2;Q $\to$ 1,2;R $\to$ 1,4,5;S $\to$ 1,2}
\opt{B}{P $\to$ 2;Q $\to$ 1,2,3,4,5;R $\to$ 1;S $\to$ 4}
\opt{C}{P $\to$ 2;Q $\to$ 1,2,3,4,5;R $\to$ 2,5;S $\to$ 1,4}
\opt{D}{P $\to$ 3;Q $\to$ 2;R $\to$ 4;S $\to$ 4}
\end{qoptions}
\sect{SECTION-II}
\qn{1} Given $\displaystyle\sum_{n=1}^\infty\dfrac1{n^2}=\dfrac{\pi^2}6$ and $\displaystyle\int_0^1\dfrac{\log_ex}{1-x^2}dx=-\dfrac{\pi^2}\lambda$, then the value of $\lambda$ equals
\qn{2} The value of $\displaystyle\lim_{n\to\infty}\dfrac1n\sum_{k=1}^n\left[e^{(1+k/n)^2}-\dfrac{3e^{\left(1+\frac{3k}n\right)}}{2\sqrt{1+\frac{3k}n}}\right]$ equals
\qn{3} Let $f(x)=\dfrac x{1+(\ell nx)(\ell nx)(\ell nx)\dots\infty}$ $\forall x\in[1,3]$, is non-differentiable at $x=k$. Then the value of $[k^2]$, is ($[.]$ denotes the greatest integer function)
\qn{4} Let $C_1$, $C_2$ and $C_3$ be the three circles lying in first quadrant each touching both axes. $C_1$ touches $C_2$ and $C_2$ touches $C_3$. If area of smallest circle is 1 unit and area of largest circle is $\left(a+2\sqrt b\right)^4$, (where $a$ and $b$ are coprime) then value of $a+b$ is
\qn{5} A circle of radius $2\sqrt{10}$ cm is divided by chord of length 12cm into two parts. If maximum area of rectangle inscribed into smaller circular segment is $\lambda$, then $\dfrac\lambda{\sqrt{15}}$ is
\qn{6} The value of
\par\noindent\hfill{\small$\displaystyle\lim_{x\to\frac\pi4}\left(\dfrac{4\sqrt2-(\cos x+\sin x)^5-10\cos x+10\cos x\sin2x}{1-\sin2x}\right)$}\hfill\null\par\noindent is
\columnbreak
\papersec{2}
\partsec{1}{PHYSICS}
\sect{SECTION-I(i)}
\qn{1} In the given circuit, initial charge on the capacitor is 1.5 C with polarity as shown in figure and current in the inductor is zero. Now at $t=0$ if the key $k$ is closed, then the minimum current through the key is :-
\qfig{p2_phy_s1_q01_1}
\begin{qoptions}
\opt{A}{33/25 A}
\opt{B}{34/25 A}
\opt{C}{36/25 A}
\opt{D}{None}
\end{qoptions}
\qn{2} A coil (which can be modelled as a series circuit) has been designed for high $Q$ performance at a rated voltage and a specified frequency. If the frequency of operation is doubled and the coil is operated at the same rated voltage, then the $Q$ factor and the active power $P$ consumed by the coil will be affected as below :-
\begin{qoptions}
\opt{A}{$P$ is doubled, $Q$ is halved.}
\opt{B}{$P$ is halved, $Q$ is doubled.}
\opt{C}{$P$ remains constant but $Q$ increases 4 times.}
\opt{D}{$P$ decreases 4 times but $Q$ is doubled.}
\end{qoptions}
\qn{3} Given a ring of mass $M$, radius $R$ resting on a rough horizontal plane, with coefficient of friction `$\mu$'. A cylindrical, vertical magnetic field passing through ring is switched on at $t=0$ which varies with time as $B=Kt^2$. Ring has charge $Q$ distributed uniformly over it. Find the time after which ring starts rotating.
\begin{qoptions}
\opt{A}{$\dfrac{\mu Mg}{QKR}$}
\opt{B}{$\dfrac{2\mu Mg}{QKR}$}
\opt{C}{$\dfrac{\mu Mg}{2QKR}$}
\opt{D}{$\dfrac{\mu Mg}{4QKR}$}
\end{qoptions}
\qn{4} The circuit shown in the diagram extends to the right into infinity. Each battery has the electromotive force $V$ (unknown) and the internal resistance $r$ (known). Each resistor has the resistance $4r$. The reading of the ideal ammeter shown in the diagram is $I$. Find the value of $V$ in terms of $I$ and $r$.
\qfig{p2_phy_s1_q04_1}
\begin{qoptions}
\opt{A}{$V=\left(2-2\sqrt2\right)Ir$}
\opt{B}{$V=\left(2+2\sqrt2\right)Ir$}
\opt{C}{$V=4.5\,Ir$}
\opt{D}{$V=3.0\,Ir$}
\end{qoptions}
\sect{SECTION-I(ii)}
\qn{5} In the circuit shown, which of the following statement is/are correct?
\qfig{p2_phy_s2_q05_1}
\begin{qoptions}
\opt{A}{Power factor of the circuit is $\dfrac1{\sqrt2}$.}
\opt{B}{Reactance of the circuit is equal to resistance of the circuit.}
\opt{C}{Circuit is in a state of resonance.}
\opt{D}{Average power dissipated in the circuit is 100 W.}
\end{qoptions}
\qn{6} A charged particle having its charge to mass ratio as $\beta$ goes in a conical pendulum of length $L$ making an angle $\theta$ with vertical and angular velocity $\omega$. If a magnetic field $B$ is directed vertically downwards (see figure) :
\qfig{p2_phy_s2_q06_1}
\begin{qoptions}
\opt{A}{$B=\dfrac1\beta\left[\omega-\dfrac g{\omega L\cos\theta}\right]$}
\opt{B}{Angular momentum of the particle about the point of suspension remains constant.}
\opt{C}{If the direction of $B$ were reversed maintaining same $\omega$ and $L$, then $\theta$ will remain unchanged.}
\opt{D}{Rate of change of angular momentum of the particle about the point of suspension is not a constant vector.}
\end{qoptions}
\qn{7} Three capacitors each having capacitance $C=2\,\mu$F are connected with a battery of emf 30 V as shown in figure. Select the correct statement(s), when the switch $S$ is closed.
\qfig{p2_phy_s2_q07_1}
\begin{qoptions}
\opt{A}{The amount of charge flown through the battery is 20 $\mu$C.}
\opt{B}{The heat generated in the circuit is 0.6 mJ.}
\opt{C}{The energy supplied by the battery is 0.6 mJ.}
\opt{D}{The amount of charge flown through the switch $S$ is 60 $\mu$C.}
\end{qoptions}
\columnbreak
\sect{SECTION-II(i)}
\paraq{1}{2}
A conducting rod of negligible resistance can slide on smooth U shaped rail as shown. At $t=0$ conducting rod is given a velocity of 10 m/s towards right. A uniform magnetic field $B=1$ T exists in the region. [Mass of conducting rod is 1 kg and $\ell=1$ m]
\qfig{p2_phy_s3_para1_2_1}
\qn{1} The distance travelled (in m) by the conducting rod when its speed is 5 m/s is :
\qn{2} Heat loss (in J) in the resistor during the time interval $t=0$ to time $t$ when the speed of the conducting rod is 6 m/s is :
\paraq{3}{4}
An infinitely long wire lying along $z$-axis carries a current $I$, flowing towards positive $z$-direction. There is no other current, consider a circle in $x$-$y$ plane with centre at (2 meter, 0, 0) and radius 1 meter. Divide the circle in small segments and let $d\vec\ell$ denote the length of a small segment in anticlockwise direction, as shown.
\qfig{p2_phy_s3_para3_4_1}
\qn{3} The maximum value of path integral $\displaystyle\int\vec B\cdot d\vec\ell$ of the total magnetic field $\vec B$ along the perimeter of the given circle between any two points on the circle is found to be $\dfrac{\mu_0I}N$, find $N$.
\qn{4} Consider two points $A(3,0,0)$ and $B(2,1,0)$ on the given circle. The path integral $\displaystyle\int_A^B\vec B\cdot d\vec\ell$ of the total magnetic field $\vec B$ along the perimeter of the given circle from $A$ to $B$ is found to be $\dfrac{\mu_0I}{2\pi}\tan^{-1}\alpha$, find $\alpha$.
\sect{SECTION-II(ii)}
\qn{5} Two metal balls of radii $R_1$ and $R_2$ such that $R_2=2R_1$ is connected to a long thin conducting wire and has total charge $Q=21\mu$C. Then the ball of radius $R_1$ is placed inside a grounded metal sphere of radius $R=3R_1$ as shown. What charge (in $\mu$C) will flow through the wire in this process ?
\qfig{p2_phy_s4_q05_1}
\qn{6} A ring of radius $R$ having a linear charge density $\lambda$ moves towards a solid imaginary sphere of radius $\dfrac R2$, so that the centre of ring passes through the centre of sphere. The axis of the ring is perpendicular to the line joining the centres of the ring and the sphere. If the maximum flux through the sphere in this process is $\dfrac{2\lambda\pi R}{n\in_0}$ then find the value of $n$?
\qfig{p2_phy_s4_q06_1}
\qn{7} A metal is heated in a furnace where a sensor is kept above the metal surface to read the power radiated ($P$) by the metal. The sensor has a scale that displays $\log_2(P/P_0)$, where $P_0$ is a constant. When the metal surface is at a temperature of $487^\circ$C, the sensor shows a value 1. Assume that the emissivity of the metallic surface remains constant. What is the value displayed by the sensor when the temperature of the metal surface is raised to $2767^\circ$C ?
\qn{8} A rod of length $\ell=2$ and varying cross section as shown is acted upon by two equal and opposite forces $F=4$N. Elastic constant of material of rod is $E=2\times10^{11}$ N/m$^2$. The increase in length of the rod is $x_0\times10^{-11}\ \ell$n(2)m. Find $x_0$. [Given, $b=1$m, $h_1=1$m and $h_2=2$m]
\qfig{p2_phy_s4_q08_1}
\qn{9} A charge $q_0$ is free to move in a circular groove placed in a solenoid as shown. Charge particle is released from position `$A$' in the groove. The speed of particle when it reaches `$B$', is proportional to $r^n$. Given that current in solenoid increases at rate of $K$ Amp/sec. Then value of $n$ is :
\qfig{p2_phy_s4_q09_1}
\qn{10} Two electrically insulated coils are wound on the same iron core. The self inductance of coils are 0.01 H and 0.04 H. The resistance of second coil is 0.3$\Omega$. What can be the maximum current (in A) in the second coil if the current in 1$^{\text{st}}$ coil is changed uniformly at a rate of 30 A/s.
\partsec{2}{CHEMISTRY}
\sect{SECTION-I(i)}
\qn{1} Which of the following statement is correct?
\begin{qoptions}
\opt{A}{In Mond's process a complex is formed by metal involves d-orbital having zero nodal plane in hybridisation of central atom}
\opt{B}{In refining of gold impurity of Cu is removed in the form of copper metaborate which is highly soluble in water.}
\opt{C}{Tempering is process in which steel becomes quite hard and brittle.}
\opt{D}{In electrorefining of Aluminium molten mixture is used as electrolyte in which Cryolite is present but Fluorspar is not present.}
\end{qoptions}
\qn{2} The crystal field splitting, $\Delta_o$ depends upon the field produced by the ligand and charge on the metal ion. Some ligands are able to produce strong fields in which case, the splitting will be large whereas others produce weak field and consequently result in small splitting of d-orbitals.

\par\noindent\hfill$\begin{array}{c}\mathrm{(A)}\quad\mathrm{(B)}\\\mathrm{[Co(NH_3)_6]^{+2}+en\rightarrow[Co(NH_3)_4(en)]^{+2}}\\\Big\downarrow\ \mathrm{en}\\\mathrm{[Co(NH_3)_2(en)_2]^{+2}\ (C)}\\\Big\downarrow\ \mathrm{en}\\\mathrm{[Co(en)_3]^{+2}\ (D)}\end{array}$\hfill\null\par
Which of the following statement is correct regarding above process?
\begin{qoptions}
\opt{A}{Total number of geometrical isomer in complex C is 2 and total optical isomer is 3.}
\opt{B}{Maximum number of atom which can lie in same plane in complex B is 7}
\opt{C}{Sum of total number of optical active isomer in complex B and number of unpaired electron in complex A is 3}
\opt{D}{Complex A has $\mu=1.732$ B.M. and complex D does not show geometrical isomer.}
\end{qoptions}
\qn{3} \qfigi{p2_chem_s1_q03_1}$\xrightarrow{\mathrm{KSH}}$ product of the reaction is?
\begin{qoptions}
\opt{A}{Ph$-\overset{\star}{\mathrm{C}}$H$_2-$CH$_2-$SH only}
\opt{B}{Ph$-$CH$_2-\overset{\star}{\mathrm{C}}$H$_2-$SH only}
\opt{C}{Both (A) and (B)}
\opt{D}{None of these}
\end{qoptions}
\qn{4} Which of the following explains the increase of reaction rate by a catalyst ?
\begin{qoptions}
\opt{A}{Catalyst provides the necessary energy to the colliding molecules to cross the barrier.}
\opt{B}{Catalyst decreases the rate of backward reaction so that the rate of forward reaction increases.}
\opt{C}{Catalyst decreases the enthalpy change of the reaction.}
\opt{D}{Catalyst provides an alternative path of lower activation energy.}
\end{qoptions}
\sect{SECTION-I(ii)}
\qn{5} $\begin{array}{l}\mathrm{CH_2-OH}\\|\\\mathrm{CH-OH}\\|\\\mathrm{CH-OH}\\|\\\mathrm{CH_2-OH}\end{array}\ \xrightarrow[\text{excess}]{\mathrm{HIO_4}}$

which of the following can be formed by above reaction?
\begin{qoptions}
\opt{A}{$\mathrm{H-\overset{O}{\overset{\|}{C}}-H}$}
\opt{B}{CH$_3-$OH}
\opt{C}{CO$_2$}
\opt{D}{$\mathrm{H-\overset{O}{\overset{\|}{C}}-OH}$}
\end{qoptions}
\qn{6} Which of the following reaction are correctly matched?
\begin{qoptions}
\opt{A}{$\begin{array}{c}\mathrm{CH_3}\\|\\\mathrm{CH_3-C-O-CH_3}\\|\\\mathrm{CH_3}\end{array}\xrightarrow{\mathrm{HI}}\mathrm{S_N1}$}
\opt{B}{$\begin{array}{c}\mathrm{CH_3-CH-O-CH_3}\\|\\\mathrm{CH_3}\end{array}\xrightarrow{\mathrm{H_2O/H^+}}\mathrm{S_N1}$}
\opt{C}{$\begin{array}{c}\mathrm{CH_3-CH-O-CH_3}\\|\\\mathrm{CH_3}\end{array}\xrightarrow{\mathrm{HI}}\mathrm{S_N2}$}
\opt{D}{$\mathrm{CH_3{-}O{-}CH_2-CH_3}\xrightarrow{\mathrm{HBr}}\mathrm{S_N2}$}
\end{qoptions}
\qn{7} Which of the following expressions are correct for an ideal gas :
\begin{qoptions}
\opt{A}{$\left(\dfrac{\partial U}{\partial V}\right)_T=0$}
\opt{B}{$\left(\dfrac{\partial U}{\partial P}\right)_T=0$}
\opt{C}{$\left(\dfrac{\partial S}{\partial V}\right)_T=0$}
\opt{D}{$\left(\dfrac{\partial H}{\partial V}\right)_T=0$}
\end{qoptions}
\qfullbegin
\sect{SECTION-II(i)}
\paraq{1}{2}
Qualitative analysis of inorganic salts the identification of cations and anions present in the salt or a mixture of salts. In the formation of a salt, the part contributed by the acid is known as anion and the part contributed by the base is known as cation.
\qfig{p2_chem_s3_para1_2_1}
\qfullend
\qn{1} Find value of $(x+2y+3z)$.

where

$x=$ Number of directional orbital used by central atom in hybridisation of anionic part of (G).

$y=$ Total number of geometrical isomers of compound (G).

$z=$ Total number of stereoisomers and unpaired electrons in compound (G).
\qn{2} Find sum of oxidation state of central atom of anion in compound (H) and number of atoms in anion of compound (I) which are in zero oxidation state.
\paraq{3}{4}
The solubility of any solid in water may be classified as

Molecular solubility : $\mathrm{AB(s)}\overset{\mathrm{aq.}}{\rightleftharpoons}\mathrm{AB(aq)}$

Ionic solubility : $\mathrm{AB(s)}\overset{\mathrm{aq.}}{\rightleftharpoons}\mathrm{A^+(aq)+B^-(aq)}$

For a substance AB(s) the following thermodynamics informations are available at 300 K :

Molecular solubility : $\Delta H^\circ=+20$ kcal/mol; $\Delta S^\circ=+40$ cal/K-mole

Ionic solubility : $\Delta H^\circ=-25$ kcal/mol; $\Delta S^\circ=-50$ cal/K-mole.
\qn{3} How many of the following statement(s) are \textbf{incorrect} for the solubility of AB(s) in water at 300 K :

(1) It undergoes molecular solubility

(2) It undergoes ionic solubility

(3) It undergoes ionic as well as molecular solubility

(4) It is completely insoluble in water.
\qn{4} Calculate the value of $\ln K_{eq}$ at 300 K for the reaction

$\mathrm{AB(aq)\rightleftharpoons A^+(aq)+B^-(aq)}$
\sect{SECTION-II(ii)}
\qn{5} How many statements is/are correct?

(i) Both PO$_4^{-3}$ \& AsO$_4^{-3}$ give canary yellow ppt. with Ammonium Molybdate under hot \& normal temperature condition respectively.

(ii) Dilute salt of chloride also give chromyl chloride test and it is non-redox reaction

(iii) NO$_3^-$ gives brown ring complex with H$_2$SO$_4$ irrespective of concentration of H$_2$SO$_4$.

(iv) Dichromate ion is more stable in acidic medium and chromate ion is more stable in basic medium
\qn{6} How many following statements is/are correct?

(i) Except alkali metal carbonate, all carbonates are insoluble in water.

(ii) Cl$_2$ water and H$_2$O$_2$ solution both are used to differentiate CO$_2$ and SO$_2$ gas.

(iii) Iodine Azide test is specific test of Sulphide ion.

(iv) Na$_2$SO$_3$ and NaHSO$_3$ can be differentiated by Litmus paper.
\qn{7} \qfig{p2_chem_s4_q07_1}
Number of $-$OH groups in (P) $=$ ($x$)

Number of $-$OH group in (Q) $=$ ($y$)

Number of alcohol group in (R) $=$ ($z$)

Number of ketone groups in (S) $=$ ($k$)

the value of $x+y+z+k$ is?
\qn{8} How many of the following gives Iodoform reaction?
\qfig{p2_chem_s4_q08_1}
\qfig{p2_chem_s4_q08_2}
\qfig{p2_chem_s4_q08_3}
\qfig{p2_chem_s4_q08_4}
\qfig{p2_chem_s4_q08_5}
\qn{9} Find freezing point temperature (in Kelvin)of liquid if vapour pressure of solid \& liquid are given by following expression.

$\ell np_s=10-\dfrac{3000(\mathrm K)}T$ where $P_s=$ vapour pressure of solid

$\ell np_\ell=5-\dfrac{2000(\mathrm K)}T$ \quad $P_\ell=$ vapour pressure of liquid

\textbf{Note :} All terms are in SI unit
\qn{10} The mole fraction of a solute in a solution is 0.1. At 298 K, molarity of this solution is the same as its molality. Density of this solution at 298 K is 2.0 g cm$^{-3}$. The ratio of the molecular weights of the solute and solvent, $\left(\dfrac{MW_{\text{solute}}}{MW_{\text{solvent}}}\right)$, is
\partsec{3}{MATHEMATICS}
\sect{SECTION-I(i)}
\qn{1} The absolute value of the expression $y^3\dfrac{d^2y}{dx^2}$ for the ellipse $3x^2+4y^2=12$ is $\dfrac PQ$ (where $P$ and $Q$ are coprime) then the value of $(P-Q)$ is
\begin{qoptions}
\opt{A}{2}
\opt{B}{3}
\opt{C}{4}
\opt{D}{5}
\end{qoptions}
\qn{2} Let $f(x)=\displaystyle\lim_{h\to0}\dfrac{(\sin(x+h))^{\ell n(x+h)}-(\sin x)^{\ell nx}}h$ then $f\left(\dfrac\pi2\right)$ is
\begin{qoptions}
\opt{A}{0}
\opt{B}{1}
\opt{C}{$\ell n\dfrac\pi2$}
\opt{D}{Does not exist}
\end{qoptions}
\qn{3} A circle touches two perpendicular lines $OA$ and $OB$ at points $M$ and $N$ respectively and passes through the vertex $C$ of rectangle $OACB$ (in this order). If the perpendicular distance of $C$ from the line $MN$ is of 10 units length, then the area of the rectangle $OACB$ is
\begin{qoptions}
\opt{A}{50}
\opt{B}{100}
\opt{C}{225}
\opt{D}{25}
\end{qoptions}
\qn{4} If $C_1\equiv y=\dfrac1{1+x^2}$ and $C_2\equiv y=\dfrac{x^2}2$ be two curve lying in $XY$ plane. Then
\begin{qoptions}
\opt{A}{area bounded by curve $y=\dfrac1{1+x^2}$ and $y=0$ is $2\pi$}
\opt{B}{area bounded by $C_1$ and $C_2$ is $\dfrac\pi2-\dfrac13$}
\opt{C}{area bounded by $C_1$ and $C_2$ is $1-\dfrac\pi2$}
\opt{D}{area bounded by curve $y=\dfrac1{1+x^2}$ and $x$-axis is $\dfrac\pi2$}
\end{qoptions}
\sect{SECTION-I(ii)}
\qn{5} Let a differentiable function $f:R\to R$ be such that for all $x$ and $y$ in $R$ we have $2|f(x)-f(y)|\leq|x-y|$ and $f'(x)\geqslant\dfrac12$. Then number of points of intersection of the graph of $y=f(x)$ with
\begin{qoptions}
\opt{A}{the line $y=x$ is one}
\opt{B}{the curve $y=-x^3$ is one}
\opt{C}{the curve $2y=|x|$ can be zero}
\opt{D}{the curve $y^2=-x$ may be more than one}
\end{qoptions}
\qn{6} If \par\noindent\hfill$\displaystyle\begin{aligned}p=\lim_{n\to\infty}n^{-n^2}\Bigl\{&(n+2^0)(n+2^{-1})\\&(n+2^{-2})\dots\dots(n+2^{-n+1})\Bigr\}^n\end{aligned}$\hfill\null\par\noindent then which of the following is true ?
\begin{qoptions}
\opt{A}{$p$ is irrational}
\opt{B}{integer nearest to $p$ is 7}
\opt{C}{$p>\displaystyle\lim_{x\to\infty}\left(\dfrac{1+x}x\right)^{\ell nx-x}$}
\opt{D}{None of these}
\end{qoptions}
\qn{7} Let $g:R\to R$ be a function defined by $g(x)=x^3+x+4$ and $f$ is a function such that $f(g^{-1}(x))=x^3+1$, then
\begin{qoptions}
\opt{A}{$\left.\dfrac d{dx}(f(x))\right|_{x=0}=48$}
\opt{B}{$\left.\dfrac d{dx}\left(f^{-1}(x)\right)\right|_{x=9}=\dfrac1{48}$}
\opt{C}{$\displaystyle\int_0^1\dfrac{\sqrt{g(x)}}{f(x)}\left(3x^2+1\right)dx=\dfrac23\tan^{-1}\left(\dfrac{6\sqrt6-8}{1+48\sqrt6}\right)$}
\opt{D}{$\displaystyle\int_0^1\dfrac{\sqrt{g(x)}}{f(x)}\left(3x^2+1\right)dx=\dfrac23\tan^{-1}\left(\dfrac2{25}\right)$}
\end{qoptions}
\sect{SECTION-II(i)}
\paraq{1}{2}
Suppose $f(x)$ and $g(x)$ are differentiable function such that $x\,g(f(x))\cdot f'(g(x))\cdot g'(x)=f(g(x))\cdot g'(f(x))\cdot f'(x)$ $\forall\ x\in R$ and $f(x)$ \& $g(x)$ are positive for all $x\in R$. Also $\displaystyle\int_0^xf(g(t))dt=\dfrac12(1-e^{-2x})$ $\forall\ x\in R$, $g(f(0))=1$ and $h(x)=\dfrac{g(f(x))}{f(g(x))}$ $\forall\ x\in R$.
\qn{1} If the graph of $y=h(x)$ is symmetric with respect to the line $x=\lambda$, then $\lambda$ equals
\qn{2} If the largest possible value of $h(x)$ $\forall\ x\in R$ is $t$, then $[t]$ equals (where $[\cdot]$ denotes greatest integer function)
\paraq{3}{4}
Let $f(x)$ be defined in the interval $(-\infty,\infty)$ such that $f(x)=\begin{cases}\dfrac{a(1-x\sin x)+b\cos x+5}{x^2},&x<0\\3,&x=0\\\left(1+\left(\dfrac{cx+dx^3}{x^2}\right)\right)^{\frac1x},&x>0\end{cases}$ and

$g(x)$ be defined in the interval $(c,-b)$ such that $g(x)=\begin{cases}\dfrac{x^2+a}{x^2+1},&c<x\leqslant2+c\\\dfrac14\left(x^3-x^2+c\right),&2+c<x\leqslant-b-1\\\dfrac94\left(|x+b|+|x+2a|\right),&-b-1<x<-b\end{cases}$

On the basis of above information, answer the following questions :
\qn{3} If $f(x)$ is continuous at $x=0$, then $\dfrac{\left|a+b+c+e^d\right|}4$ equals
\qn{4} The number of points of non-differentiability of $g(x)$ in $x\in(c,-b)$
\sect{SECTION-II(ii)}
\qn{5} Let $\lambda=\displaystyle\int_1^9\left(\dfrac12+\sqrt{\dfrac14+\log_3x}\right)dx+\int_1^2\dfrac{3^{x^2-[x]}}{3^{\{x\}}}dx$ then $\lambda$ equals

(where $[.]$ denotes greatest integer function and $\{.\}$ denotes fractional part function)
\qn{6} The parabola $y^2=x-y$ divides the triangle formed by the tangent to it at the origin and the straight line $x+y=0$ and $y=-1$ into two parts whose areas are $A_1$ and $A_2$ (Given $A_1<A_2$). If $\dfrac{A_1}{A_2}=\dfrac pq$ where $p,q\in N$ and $\dfrac pq$ is in simplest form then the value of $(p+q)$ is :
\qn{7} Let $L_1$, $L_2$ and $L_3$ be the length of tangents drawn from a point $P$ to the circles $x^2+y^2=4$, $x^2+y^2-4x=0$ and $x^2+y^2-4y=0$ respectively. If $L_1^4=L_2^2L_3^2+16$, then locus of $P$ are the curves $C_1$(a straight line) and $C_2$(a circle). Let the circumcentre of the triangle formed by $C_1$ and other two lines which are inclined at an angle of $45^\circ$ with $C_1$ and is tangent to $C_2$ is $(\alpha,\beta)$, then $(\alpha+2\beta)$ equals
\qn{8} Find the values of $x$ satisfying the equation $2\cos^{-1}(1-x)-3\cos^{-1}x=\pi$
\qn{9} Let $f(x)$ be a differentiable function on $[0,8]$ such that $f(1)=5$, $f(2)=\dfrac12$, $f(3)=6$, $f(4)=-3$, $f(5)=4$, $f(6)=\dfrac14$ and $f(7)=\dfrac{-1}5$. If the minimum number of roots of the equation $f'(x)-f'(x)(f(x))^2=0$ is $\lambda$, then $\lambda$ equals
\qn{10} Let $P(x)$ be a polynomial of degree 4 having a local maximum at $x=2$ and $\displaystyle\lim_{x\to0}\left(3-\dfrac{P(x)}x\right)=27$. If $P(1)=-9$ and $P''(x)$ has a local minimum at $x=2$. If global maximum value of $y=P'(x)$ on the set $A=\{x:x^2+12\leq7x\}$ is $4M$ then value of $M$ is :
\end{multicols}
\end{document}
